\documentclass{article}
\usepackage[T1]{fontenc}
\usepackage{lmodern}
\usepackage{graphicx}
\usepackage{amsmath,amssymb}
\usepackage[a4paper,margin=2cm]{geometry}
\usepackage[absolute,overlay]{textpos}
\setlength{\TPHorizModule}{1mm}
\setlength{\TPVertModule}{1mm}
\usepackage[indonesian]{babel}
\usepackage{hyperref}
\setlength{\emergencystretch}{2em}
% unit: Halaman 40

\begin{document}

% === Halaman 40 (python/native) ===

40

1. RC4

• Cipher alir yang paling popular, namun sekarang sudah tidak aman

• Dibuat oleh Ronald Rivest (1987) dari Laboratorium RSA

• RC adalah singkatan dari Ron’s Code. Versi lain mengatakan Rivest 

Cipher .

• Digunakan di dalam beberapa sistem keamanan seperti:

 - protokol SSL (Secure Socket Layer) dan TLS (Transport Layer Security) 

 - Standard IEEE 802.11 wireless LAN: WEP (Wired Equivalent Privacy)

 - WPA (Wi-fi Protocol Access) untuk nirkabel

\end{document}
